\chapter{Negative Binomial Distribution}
\label{chap:negative-binomial-distribution}
\lecture{5}{1 Jun. 10:18}{Fifth Lecture}

\begin{eg}[Coin Tossing for Multiple Heads]
  \label{eg:coin-toss-nbd}
  Suppose we are flipping a fair coin repeatedly. Each flip is an independent
  Bernoulli trial. Let \(p\) be the probability of getting a Head (success), and
  \(q=1-p\) be the probability of getting a Tail (failure). We are interested in
  the number of Tails, \(X\), observed before we obtain exactly \(r\) Heads.
  \marginpar{Example: Number of tails before \(r\) heads.}

  For instance, let \(r=3\) (we want 3 Heads) and assume the probability of a
  Head is \(p=0.5\). We continue flipping the coin until we have observed 3
  Heads. The number of Tails(failures) \(X\) that occurred before the 3rd Head
  is a random variable following a Negative Binomial distribution.
  Possible outcomes for \(X\) could be:
  \begin{itemize}
    \item HHH: 0 Tails (\(k=0\)). The sequence ends with the 3rd Head.
    \item THHH: 1 Tail (\(k=1\)).
    \item HTHH: 1 Tail (\(k=1\)).
    \item HHTH: 1 Tail (\(k=1\)).
    \item TTHHH: 2 Tails (\(k=2\)).
  \end{itemize}
  This scenario perfectly illustrates the conditions for a Negative Binomial
  distribution:
  \begin{itemize}
    \item The experiment consists of a sequence of independent trials (each coin
      flip).
    \item Each trial has only two possible outcomes: success (Head) or failure
      (Tail).
    \item The probability of success (\(p\)) is constant for each trial.
    \item The experiment continues until a fixed number of successes (\(r\)) is
      achieved.
  \end{itemize}
  The random variable of interest, \(X\), is the number of failures (Tails) that
  occur before the \(r\)-th success (Head).
\end{eg}

\section{Definition}
\begin{definition}[Negative Binomial Distribution]
  \label{def:neg-binomial}
  The Negative Binomial distribution, sometimes referred to as the Pascal
  distribution, is a discrete probability distribution. It models the random
  variable \(X\) that represents the \(\#\) failures that occur
  before a specified (constant or fixed) number of successes, denoted by \(r\),
  is achieved.\\

  \textbf{A random variable \(X\) that follows a Negative Binomial distribution is
  commonly denoted as \(X \sim \text{NB}(r, p)\)}\\
  where,\\
  \(r\) is the \(\#\) success\\
  \(p\) is the probability of success for one trial.\\

  A random process is described by the Negative Binomial distribution if it
  adheres to the following conditions:
  \begin{itemize}
    \item The process is a sequence of independent trials.
    \item Each trial results in one of two outcomes, termed success or failure.
    \item The probability of success, \(p\), is constant for all trials.
    \item The experiment is continued until a predetermined total of \(r\)
      successes are observed. Here, \(r\) must be a positive integer.
  \end{itemize}
\end{definition}

\begin{remark}
  \label{remark:alter-formulation-nbd}
  It's worth noting an alternative formulation where the distribution
  models the total number of trials \(Y = X+r\) required to achieve \(r\)
  successes.
\end{remark}

\section{Probability Mass Function (PMF) of \(X\)}
The \textbf{PMF} of \(\text{NB}(r,p)\), \(P(X=x)\), quantifies the probability
of observing exactly \(x\) failures prior to achieving the \(r\)-th success in a
series of independent Bernoulli trials with \(p\) the success probability.

That means, there are \(x\) failures and \(r\) success. So the total \(\#\)
trials is \(k=x+r\). And first \(x+r-1\) trials represent \(Y \sim \text{Bin}(n,p)\).
Then,
\begin{align*}
  P(X=x) &= P(Y=r-1)\cdot p\\
  \marginpar{\(\because\) probability of \(k^{th}\) success is \(p\)}\\
  P(X=x) &= \binom{x+r-1}{r-1}q^{x}p^{r-1}\cdot p\\
  P(X=x) &= \binom{x+r-1}{r-1}q^{x}p^{r}\\
\end{align*}

\(\therefore\) PMF for \(X \sim NB(r,p)\) is:

\begin{equation}
  \label{eq:nbd_pmf}
  \boxed{
    P(X=x) = \binom{x+r-1}{x} p^r (1-p)^x, \quad \text{for } x\in \mathbb{N}_0
  }
\end{equation}

Here, \(1-p\) is often denoted as \(q\), the probability of failure.

\section{Expectation of \(X\)}
\label{sec:nbd_expectation}
The expectation (or mean) of a random variable \(X \sim \text{NB}(r, p)\),
denoted \(E[X]\), can be derived by considering \(X\) as a sum of independent
Geometric random variables. 

Let \(X_i\) for \(i=1, \dots, r\) be the number of failures between the
\((i-1)\)-th success and the \(i\)-th success.

\begin{figure}[H]
  \centering
  \includegraphics[width=0.8\textwidth]{Figures/pimg2.png}
  \caption{\(X_i\) follows Geometric distribution}
  \label{fig:pimg2--x_i-follows-geometric-distribution}
\end{figure}

Each \(X_i\) follows a Geometric distribution with success probability \(p\),
specifically \(X_i \sim Geom(p)\). The PMF for such a variable (counting
failures before the first success) is \(P(X_i=k) = (1-p)^k p\) for \(k=0, 1, 2,
\dots\). The expectation of this form of Geometric distribution is \(E[X_i] =
\frac{q}{p}\). \marginpar{For \(Y \sim Geom(p)\) counting failures before 1st
success, \(E[Y] = q/p\).}\\

The total number of failures \(X\) before \(r\) successes is the sum of these
\(r\) independent Geometric variables:
\[ X = X_1 + X_2 + \dots + X_r \]
By the linearity of expectation:
\begin{align*}
  E[X] &= E[X_1 + X_2 + \dots + X_r] \\
       &= E[X_1] + E[X_2] + \dots + E[X_r] \\
       &= \sum_{i=1}^{r} E[X_i] \\
       &= \sum_{i=1}^{r} \frac{q}{p} \\
       &= r \frac{1-p}{p}\\
       &= \frac{rq}{p}
\end{align*}
Thus, the expectation of \(X \sim NB(r,p)\) is:
\begin{equation}
  \label{eq:nbd_expectation}
  E[X] = \frac{rq}{p}
\end{equation}
\(\E[X]\) is expected number of failures before r success

\section{Variance of \(X\)}
\label{sec:nbd_variance}
The variance of \(X \sim NB(r, p)\), denoted \(Var(X)\), can similarly be
derived using the representation of \(X\) as the sum of \(r\) independent
Geometric random variables \(X_i \sim Geom(p)\), as defined in the expectation
section. The variance of each \(X_i\) (number of failures before the first
success, with success probability \(p\)) is \(Var(X_i) = \frac{q}{p^2}\).
\marginpar{For \(Y \sim Geom(p)\) counting failures before 1st success, \(Var(Y)
= (q)/p^2\).}\\

Since the \(X_i\) variables are independent, the variance of their sum is the sum of their variances:
\begin{align*}
  Var(X) &= Var(X_1 + X_2 + \dots + X_r) \\
         &= Var(X_1) + Var(X_2) + \dots + Var(X_r) \quad \text{(due to independence of \(X_i\))} \\
         &= \sum_{i=1}^{r} Var(X_i) \\
         &= \sum_{i=1}^{r} \frac{q}{p^2} \\
         &= \frac{rq}{p^2}
\end{align*}
Thus, the variance of the Negative Binomial distribution (\(X \sim NB(r,p)\)) is:
\begin{equation}
  \label{eq:nbd_variance}
  Var(X) = \frac{rq}{p^2}
\end{equation}
An important property is that the variance of the Negative Binomial distribution
can be larger than its mean, making it suitable for modeling overdispersed count
data.

\section{Alternate Formulation using \(Y=X+r\)}
\label{sec:alternate-formulation-using-y-x-r-}

Here, \(Y\) is the random variable that represents the \(\#\) trials to get
\(r^{th}\) success.

\subsection{PMF of Y}
\label{ssec:pmf-of-y-nbd}

\begin{equation}
  \label{eqn:pmf-of-y-nbd}
  \boxed{P[Y=y] = \binom{y-1}{r-1}q^{y-r}\cdot p^{r}}
\end{equation}

This measures the probability of getting \(r^{th}\) success in \(y\) trials.

\subsection{Expectation of Y}
\label{ssec:expectation-of-y-nbd}

\marginpar{Using \href{ssec:linearity-of-expectation}{Linearity of Expectation}}
\begin{equation}
  \label{eqn:expectation-of-y-nbd}
  \boxed{\E[Y]=\E[X+r]=\frac{r(1-p)}{p}+r = \frac{r}{p}}
\end{equation}
This measures the expected number of trials to get \(r^{th}\) success.

\subsection{Variance of Y}
\label{ssec:variance-of-y-nbd}

\begin{equation}
  \label{eqn:variance-of-y-nbd}
  \boxed{\text{Var}[Y] = \text{Var}[X+r]= \text{Var}[X]  = \frac{rq}{p^2}}
\end{equation}
